\section{Reflection}

This project provided a comprehensive hands-on experience with every stage of a real machine
learning pipeline applied to physiological signal data. The most technically demanding aspect
was implementing the entropy feature extraction --- particularly Sample Entropy and LZiv
Complexity --- which required careful implementation using explicit \texttt{float64} conversion
and binary signal thresholding to ensure numerical stability across all 5\,850 trials.
Debugging the Kruskal-Wallis \texttt{ValueError} caused by zero-variance features after
standardisation was another instructive challenge, resolved by wrapping the test in a
\texttt{try-except} block.

A key observation was the impact of feature selection on downstream classification performance.
The Kruskal-Wallis test revealed that not all 26 extracted features are statistically significant
for this classification task, and reducing the feature space before PCA led to more numerically
stable principal components. The feature-domain ablation experiment was a particularly valuable
addition to the pipeline: it transformed a broad task requirement into a concrete, quantified
experimental finding --- demonstrating that entropy features are the most informative single
domain and that frequency-domain features contribute the least when used in isolation.

Working individually rather than in a group required a more systematic approach to comparative
analysis: multiple classifiers, validation strategies, and feature sets were compared
systematically, providing analytical breadth equivalent to a group project. The addition of a
feature importance analysis and a GridSearchCV heatmap gave further insight into which features
and hyperparameters drive performance, reinforcing the experimental findings.

This project has strengthened understanding of the full ML pipeline --- from signal-level
preprocessing decisions through to the interpretation of confusion matrices, ROC curves,
cross-validation stability, and feature importance rankings. Future directions include extending
to multiclass gesture classification, incorporating multi-channel EMG, exploring subject-independent
generalisation, and investigating whether deep learning methods such as 1D CNNs or BiLSTMs can
further improve accuracy on this dataset.
